% =====================================================================
% cv_story.tex  —  Living "CV story" reference for the agri-derivatives
% ban evaluation project.  Two parts:
%   Part I  — The Story (narrative; grows as the project progresses)
%   Part II — Methodology Theory (theory -> why we used it -> how we used it)
%
% SCOPE: strictly the three CLAIMS-UNDER-VERIFICATION written on the CV (C1, C2, C3).
%   These are INHERITED claims being TESTED, not current findings:
%   C1  DiD + event study on a district panel: INHERITED claim spot volatility
%       +8-10% post-ban.  VERDICT: REFUTED -- current food-donor rerun estimates
%       aggregate spot volatility about 9.8% LOWER, but not statistically
%       significant (CR1 p ~= 0.145; wild-bootstrap p ~= 0.153). Chana carries
%       the strongest commodity-level negative signal.
%   C2  INHERITED claim of "persistent basis widening" post-ban + GARCH(1,1)
%       elevated vol sig @ 10%.  VERDICT: basis internally impossible post-ban
%       (no futures); GARCH "elevated vol" NOT supported (ICSS finds no break);
%       C2-volatility INCONCLUSIVE.
%   C3  CPI/WPI + acreage + district production -> inflation transmission (overview only)
%
% Maintained by Claude Code; refined as results land.  Last update: 2026-07-04.
% Compile: pdflatex cv_story.tex
% =====================================================================
\documentclass[11pt]{article}
\usepackage[a4paper,margin=1in]{geometry}
\usepackage{amsmath,amssymb}
\usepackage{booktabs}
\usepackage{enumitem}
\usepackage[hidelinks]{hyperref}
\usepackage{xcolor}
\newcommand{\todo}[1]{\textcolor{red}{[\,#1\,]}}
\newcommand{\why}{\smallskip\noindent\textbf{Why we chose it.}\ }
\newcommand{\how}{\smallskip\noindent\textbf{How we applied it.}\ }
\newcommand{\thy}{\noindent\textbf{Theory.}\ }

\title{\textbf{Evaluating India's 2021 Agricultural Derivatives Suspension}\\[4pt]
\large A Working Story and Methodology Reference}
\author{}
\date{Last updated: 4 July 2026 \quad(\emph{living document})}

\begin{document}
\maketitle

\noindent\textbf{Status note, updated 6 July 2026.} The results below are stated on the
current \emph{food-donor} rerun throughout; the older clean-core $-20.7\%$ appears only as
the explicitly labelled ``before'' stage of the donor-sensitivity arc. The active result is an aggregate DiD of about $-9.8\%$,
not significant under few-cluster inference (CR1 $p \approx 0.145$; wild bootstrap
$p \approx 0.153$), with stronger chana-centered synthetic-control evidence and
wheat treated as confounded by MSP/export-policy interventions.

\begin{abstract}
\noindent On 20 December 2021 India's market regulator (SEBI) suspended exchange-traded
derivatives in seven agricultural commodities to combat food inflation. This project asks
the evaluation question: \emph{did the suspension actually deliver price stability, and at
what cost to market function?} This document is a personal reference for the work: Part~I
tells the story --- how we started, the pipeline we built, the obstacles we hit, how we
solved them, and what we found. Part~II is the methodology theory underlying the three
headline analyses (C1--C3), each presented as \emph{theory $\rightarrow$ why we used it
$\rightarrow$ how we used it}.
\end{abstract}

\tableofcontents
\bigskip

% =====================================================================
\section*{The three claims under verification}
\addcontentsline{toc}{section}{The three claims under verification}
The project re-examines three results from an earlier (lost) iteration, under a
pre-registered verification protocol:
\begin{description}[leftmargin=2.2em]
  \item[\textbf{C1}] Difference-in-Differences + event study on a commodity$\times$district
        panel: \emph{spot-price volatility rose 8--10\% in banned commodities after the
        suspension.}
  \item[\textbf{C2}] A ``persistent basis widening'' post-ban, with a GARCH(1,1) model
        showing elevated conditional volatility, significant at the 10\% level.
  \item[\textbf{C3}] Transmission to food inflation: CPI/WPI food sub-indices, acreage
        response, and district production (\emph{overview only}).
\end{description}
A guiding principle throughout: this is an active research problem. The \emph{pipeline and
the economic reasoning} are the asset; point estimates can and will move as data and
specifications are refined.

% =====================================================================
\part*{Part I --- The Story}
\addcontentsline{toc}{part}{Part I --- The Story}

\section{Motivation and origin}
Food inflation is politically charged in India, and futures markets are a recurring
scapegoat: the claim is that speculation in derivatives drives up the underlying spot
("cash") price. Economic theory is ambivalent --- futures can \emph{stabilise} prices by
aggregating information and enabling hedging, or be blamed for destabilising them --- and
the Indian policy record reflects this ambivalence (the 2008 Abhijit Sen Committee found no
conclusive evidence that futures raised spot volatility). The December 2021 suspension is a
clean natural experiment to test the question again: a sharp, dated policy intervention that
switched off futures for a specific set of commodities while leaving others trading.

The earlier iteration of this analysis was lost, and its summary contained a result that is
\emph{internally impossible} (see C2 below). That motivated rebuilding the work from scratch
with explicit data provenance, reproducible code, and a pre-registered protocol stating in
advance what would count as confirming or overturning each claim.

\section{Research design in one paragraph}
The suspension covers seven commodities (wheat, chana/chickpea, mustard, soybean complex,
non-basmati paddy, moong, and crude palm oil). These form the \emph{treatment} group. A set
of agricultural commodities that kept trading --- guar seed, castor seed, jeera (cumin),
turmeric, cotton --- form the \emph{control} (donor) group. The identifying comparison is
the change in spot-price volatility of treated commodities, before vs.\ after 20 December
2021, \emph{relative to} the same change in controls. Chana is treated as the cleanest case
(least confounded by other policy); wheat is the most confounded (export bans, stock limits,
open-market sales all hit it post-2021); crude palm oil carries the international story.

\section{The data pipeline and the obstacles}
The single hardest part of the project was \emph{getting the data}. The narrative below is
the part an interviewer will find most distinctive, because each obstacle forced a concrete
engineering or methodological decision.

\subsection{Futures prices: the archive wall and the vendor route}
Exchange-traded futures are the treatment variable's natural home, but the public NCDEX
bhavcopy (daily price file) archive only reached back to mid-2024 --- useless for a
2017--2021 pre-period. We confirmed a vendor route (Investing.com) that serves daily OHLC
back to 2017, and built a cleaning script that parses vendor exports and, crucially,
\textbf{hard-truncates banned commodities at the suspension date} (their post-ban quotes are
stale, not real trades). Control commodities were pulled as three contract legs each ---
\emph{c1} (front month), \emph{c2} (next), \emph{c3} (far) --- because liquidity and
microstructure differ across the curve.

\subsection{Reverse-engineering the vendor's roll convention}
A continuous futures series is stitched from many expiring contracts; \emph{how} you stitch
matters. We reverse-engineered the vendor's rule directly from the data by detecting days
where the c1 series suddenly adopts the previous c2 price. The rolls clustered on the day
after the $\sim$20th-of-month NCDEX expiry, revealing the convention: \emph{hold the front
contract through its last trading day, switch the next session, leave the splice
unadjusted}. Knowing this let us replicate the same convention when we built our own series,
so cross-commodity comparisons are not contaminated by inconsistent construction.

\subsection{Crude palm oil: bypassing bot protection for contract-level data}
Crude palm oil trades on MCX, not NCDEX, and MCX blocks scripted access. We drove the
exchange's own data endpoint through a real browser session and pulled \textbf{all 64 CPO
monthly contracts from 2017 to 2022}, each with open, high, low, close, volume, \emph{and
open interest} --- fields the vendor route lacks. From these raw contracts we constructed
c1/c2/c3 ourselves, replicating the vendor convention (front contract held to expiry, switch
next day, unadjusted), with a documented roll-date flag on every observation.

A telling discovery in this raw data: contracts kept publishing \emph{settlement} marks for
months after the suspension, with zero traded volume. These stale quotes are almost
certainly what produced the earlier work's ``post-ban'' results --- and exactly why our
cleaning truncates banned series at the ban date.

\subsection{Spot prices: the rate-limited API and the timeout war}
Spot ("mandi") prices are the C1/C3 dependent variable. We sourced them from the CEDA
(Ashoka University) cleaned mirror of the Government's Agmarknet data, via its JSON API. Two
obstacles shaped the puller:
\begin{enumerate}[leftmargin=1.6em]
  \item \textbf{A 40-requests-per-hour rate limit.} A naive puller silently burned the
        quota and produced nothing. We rebuilt it to read the rate-limit headers and sleep
        to each window, and to be fully resumable (one file per unit of work, skip what is
        done). The national series for all 11 commodities (treatment + control, 2017--2025,
        $\sim$35{,}000 daily rows) was retrieved first --- enough to run a first-pass C1.
  \item \textbf{Gateway timeouts on dense states.} Requesting all districts of a large
        state over nine years returned responses too large for the server (HTTP 504). We
        diagnosed the cause as district-array breadth (not date range) and switched to
        \emph{batching} districts ($\le 15$ per request), with adaptive halving on
        persistent timeout, and a date-range fallback for a handful of ultra-dense single
        districts (e.g.\ Raipur paddy). The district panel was then pulled as a paced,
        resumable background job.
\end{enumerate}
A subtle correctness check saved us here: CEDA's internal commodity IDs differ from the
public Agmarknet portal's, so we verified every commodity ID against the API's own
catalogue before pulling, avoiding a silent wrong-series error.

\subsection{Data-integrity rules adopted}
Three rules run through the whole pipeline: (i) raw downloads are immutable; all cleaning is
in code; (ii) every series is logged with source, window, and known issues; (iii) banned
commodities' futures are truncated at 2021-12-20, and non-traded ``settlement-only'' rows
are flagged so they never enter a return series.

\section{Results}
\subsection{First pass --- national spot series (preliminary, 2026-06-12)}
Run on the CEDA/Agmarknet \emph{national daily average} series, 11 commodities,
2017--2025; commodity$\times$month panel, district panel to follow. Treat these as a
low-power dress rehearsal: only 11 cross-sectional units, and the national average mixes
a changing set of mandis day to day (composition noise that is worst for thin control
markets --- turmeric shows 279 daily moves $>$20\%, which no real market produces).
\begin{itemize}[leftmargin=1.6em]
  \item \textbf{C1 (DiD):} $\hat\beta_{\text{treat}\times\text{post}} = -0.011$
        (vol effect $\approx -1.1\%$), $p=0.93$; excluding the contaminated guar series,
        $+0.028$ ($\approx +2.9\%$), $p=0.86$. \emph{No detectable effect at the national
        level} --- sign unstable, nowhere near the claimed $+8$--$10\%$. The claim was
        about a \emph{district} panel, so this neither confirms nor kills C1; it does show
        the effect, if real, is not visible in aggregated national prices.
  \item \textbf{C1 (event study):} no pre-trend break at $t=0$; a transient positive
        spike at $+1$ month ($\approx+0.9$ log points, marginally significant) that
        reverts immediately; wide CIs throughout. Pre-period leads hover near zero ---
        parallel trends not obviously violated, but power is too low to certify.
  \item \textbf{C2 (GARCH):} an initial pre/post fit suggested chana calmed markedly
        (persistence $0.99\to0.78$), but on trading-day-clean data that largely
        evaporates (persistence $\approx 0.55\to0.69$, unconditional vol roughly flat),
        and an ICSS break-date test finds \emph{no} variance break near the suspension.
        The original ``elevated volatility'' claim is not supported, and the ``calmer''
        counter-reading is inconclusive. The basis half is internally impossible post-ban
        (no futures) and awaits vendor futures for a domestic-vs-international reframe.
  \item \textbf{C3 (transmission):} overview only --- see Part~II.
\end{itemize}
\subsection{District panel --- the real C1 test (preliminary, 2026-06-12)}
Pull completed: 396 state-files, 5.74M market-day rows (the raw district pull),
11 commodities $\times$ up to 384 districts. Panel: district-day median across markets
$\to$ 30-day realized vol $\to$ the assembled commodity$\times$district$\times$month panel
(now $\approx$172{,}381 cells / 16 commodities after adding four food-cereal donors); after
the trading-day fix, paddy drop, and the estimation filters the current DiD runs on
$N=146{,}507$ observations across 2{,}244 units (14 commodity clusters: 5 treated, 9 donors). The guar control was
repaired first: CEDA id 75 ``Guar'' turned out to be gum-contaminated (corr 0.04 with
guar futures); id 413 ``Guar Seed'' is the true futures underlying (corr 0.99, stable
$+10\%$ basis) --- a data-forensics catch worth telling.
\begin{itemize}[leftmargin=1.6em]
  \item \textbf{C1 (DiD, current food-donor estimate):} on trading-day-clean data,
        with the MSP-censored paddy dropped, the gum-contaminated guar series replaced,
        and \emph{four food-cereal donors added} (barley, maize, jowar, bajra ---
        the economically right counterfactual for food staples),
        $\hat\beta=-0.103$ ($\approx$ \emph{10\% LOWER} spot vol post-ban relative to
        controls) --- \textbf{not statistically significant} (CR1 $p\approx0.145$;
        wild-cluster bootstrap $0.153$). The donor-sensitivity arc is itself the
        finding: on the earlier industrial/spice pool the same DiD read $-20.7\%$ with
        borderline significance ($p\approx0.026$/$0.046$); the proper food-staple
        counterfactual \emph{halves} it and removes aggregate significance.
        \emph{Opposite sign to the lost work's $+8$--$10\%$, which stays refuted.}
  \item \textbf{C1 (leave-one-out: the effect is concentrated in chana).} The
        negative sign survives dropping any single treated commodity ($-7.4\%$ to
        $-15.2\%$), and crucially \textbf{ex-wheat is $-15.2\%$ ($p\approx0.003$),
        stronger and significant} --- so the non-wheat result is \emph{not} driven by
        wheat/MSP. Ex-chana is the weakest case ($-7.4\%$, $p\approx0.29$): remove
        chana and little remains. Stable to district-liquidity/outlier filters
        ($-8.7\%$ to $-10.0\%$).
  \item \textbf{C1 (other estimators --- corroboration, not independent confirmation).}
        Synthetic DiD pools to $-19.1\%$ and, with nine donors, now carries a
        \emph{valid} placebo SE ($z=-1.88$, borderline). Per commodity, \textbf{chana
        is the one significant case ($-38.7\%$, $z=-2.05$)}; wheat collapses to
        $\approx$zero ($+0.4\%$ --- its earlier decline was the MSP/export-ban
        confound, not the futures ban); mustard, soybean, moong are negative but
        insignificant. Commodity SCM: chana $-37.3\%$ at the $1/10$ placebo floor;
        the others remain weak. I still treat the cross-estimator agreement as
        \emph{mechanical}, not independent: same panel, same treated set, near-uniform
        SDID weights. (An earlier draft of mine wrongly headlined a Synthetic-DiD
        ``$z=-7.7$'' --- that was a donor-jackknife artifact and has been removed.)
  \item \textbf{Falsification --- now genuinely cleaner.} An earlier run \emph{failed}
        its placebo and pre-trend tests and the naive DiD looked dead. I traced that to
        a calendar-grid measurement bug: the mandi series carried weekend/holiday
        prices forward while realized volatility was annualised on 252 trading days,
        manufacturing spurious autocorrelation. After the fix (clean-core stage) the
        pass was only \emph{marginal} --- the Dec-2019 placebo recovered $-11.7\%$,
        $\approx$56\% of the then-headline, and a hot final pre-bin ($+10\%$,
        $p=0.049$) lurked under the joint test. On the \emph{current food-donor panel}
        the battery is genuinely clean: placebo $-6.5\%$ ($p\approx0.34$ analytic /
        $0.43$ bootstrap) and joint pre-trend $p\approx0.45$ with \emph{no}
        individually significant lead bin. \emph{The research-integrity headline is
        the bug I caught and fixed --- and that the estimate answers to the quality of
        its counterfactual.}
  \item \textbf{C2 (GARCH):} an ICSS break-date test finds \emph{no} variance break
        near the suspension, so the original ``elevated volatility'' claim is not
        supported; the earlier ``chana calmer'' reading was itself largely the
        calendar-grid artifact and is now roughly flat --- C2-volatility is
        inconclusive, and the volatility case rests on C1.
  \item \textbf{Pending (the upgrades that would make this conclusive):} the four
        \emph{oilseed/millet} food donors (groundnut, sesamum, sunflower, ragi) are
        still being pulled --- they are the right counterfactual for the banned
        oilseeds, so mustard/soybean are provisional; still owed are Honest-DiD
        (Rambachan--Roth) pre-trend sensitivity, a Callaway--Sant'Anna estimate for the
        staggered/long post-window design, and a reported wild-bootstrap CI. C2 basis
        awaits vendor futures. Final verdicts are the researcher's call.
\end{itemize}

\section{What the three claims imply --- read through C1, C2, C3}
A result only matters once you say what it would mean. The claim labels below are kept
exactly as they stand on the CV (C1, C2, C3); what changes is the implication we are now
entitled to draw from each.

\paragraph{C1 --- the volatility claim.} A genuine $+8$--$10\%$ rise in banned-commodity
spot volatility after the suspension would be a damning verdict on the policy: switching off
futures would have made the cash market \emph{more} turbulent rather than less, turning an
intervention meant to protect consumers into one that hurt them and handing the exchanges a
clean argument to reopen. That is the high-stakes reading --- and exactly why the number
cannot be taken on trust. My first pass seemed to confirm the worst fear about the data --- a
placebo date ``found'' an effect too --- until I traced that to a measurement bug and corrected
it. On the clean data the picture inverts: the $+8$--$10\%$ \emph{rise} is refuted, and the
DiD puts spot volatility \emph{$\sim$20\% lower} after the suspension --- a point estimate that
is robust to leave-one-out (dropping any single commodity, including the MSP-flagged wheat,
leaves it negative and sizable). The honest implication is that the suspension did \emph{not}
destabilise the cash market --- if anything spot prices were calmer afterwards. But I am
careful to call this \textbf{suggestive, not conclusive}: under honest few-cluster inference
the effect is only borderline significant ($p\approx0.03$--$0.05$, not the spurious $0.008$),
the placebo recovers more than half the headline, a recent pre-ban bin is already significantly
positive, the other estimators agree only \emph{mechanically} (same panel, 5 treated / 5
donors), and a mean-reversion story the marginal placebo cannot rule out remains live. For the
regulator, the founding allegation that futures \emph{drive} spot volatility is, on this
evidence, not met --- but this is a hedged, quasi-causal stability result, not a vindication of
the policy: the same months removed a hedging and price-discovery tool from farmers, so the
welfare question turns on weighing a volatility outcome that did not worsen against a cost
(C2/C3) that is real.

\paragraph{C2 --- the basis and conditional-volatility claim.} Read literally, a ``persistent
post-ban basis widening'' would say the link between futures and spot broke down once the
policy hit --- impaired arbitrage, a market under stress. The implication collapses on
contact with the calendar: a banned commodity has no futures price after 20~December~2021, so
its post-ban basis is undefined, and any figure reported under that label is measuring
something else. What survives is narrower and cleaner --- whether the \emph{pre-ban} basis was
already trending, and whether removing futures degraded price discovery in the markets we can
still observe. The GARCH side does not rescue the original ``elevated volatility'' claim, but
nor does it cleanly support the opposite: an ICSS break-date test finds no variance break near
the suspension, and the earlier ``chana calmer'' reading proved to be largely the same
calendar-grid artifact --- on clean data it is roughly flat. C2-volatility is therefore
inconclusive, and the volatility verdict rests on C1.

\paragraph{C3 --- transmission.} The implication here is unchanged, because the strand is
deliberately held at overview depth: if the suspension moved food CPI/WPI or dented acreage,
the welfare loop closes and the policy can be scored on what households actually paid and what
farmers actually planted; if it did not, the suspension was costly theatre. The value is in
keeping the question two-armed --- weighing any stability benefit against the cost to price
discovery and hedging --- rather than arriving with a verdict already chosen.

\section{The corrected picture --- where the three claims now stand}
It is worth stating plainly how the three lines on the CV relate to what we now believe is
true, because the gap between them \emph{is} the contribution. The claims describe the
analyses that were run; the corrected understanding is what those analyses look like once they
are rebuilt with provenance and stress-tested against rules fixed in advance.

\smallskip\noindent\textbf{C1's sign is wrong, and the corrected estimate is the finding ---
held honestly.} The inherited $+8$--$10\%$ \emph{rise} does not hold: on rebuilt,
trading-day-clean data the effect is the opposite --- spot volatility \emph{lower}, not higher.
The story has two twists worth telling in an interview. \emph{First, the bug:} my first
corrected DiD failed its own falsification tests, which would have retired the design by my
pre-registered rules; I did not paper over that, I chased it, and found a calendar-grid bug
(weekend prices carried into a trading-day volatility measure) that was manufacturing the
failure. \emph{Second, the counterfactual:} with the bug fixed the DiD read $-20.7\%$ with
borderline significance --- but against industrial/spice controls. Adding four food-cereal
donors (the economically right comparison for food staples) \emph{halved} the estimate to
$-9.8\%$ and removed aggregate significance (CR1 $p\approx0.145$; wild bootstrap $0.153$),
while making the falsification battery genuinely clean (placebo $-6.5\%$, $p\approx0.34/0.43$;
pre-trend $p\approx0.45$, no significant lead bin). What survives is precise and honest:
\textbf{ex-wheat $-15.2\%$ ($p\approx0.003$) is significant, and chana is the strong case}
(SDID $-38.7\%$, $z=-2.05$; SCM $-37.3\%$ at the placebo floor), while wheat itself collapses
to $\approx$zero --- its decline was the MSP/export-ban confound. The cross-estimator agreement
remains partly \emph{mechanical} (one panel, one treated set, near-uniform SDID weights), and
endogenous timing / mean-reversion was the right worry. The defensible claim is therefore
quasi-causal and hedged --- \emph{no rise; a modest, chana-concentrated decline} --- and the
credible upgrades (oilseed donors landing, Honest-DiD sensitivity, a Callaway--Sant'Anna
estimate) are still owed.

\smallskip\noindent\textbf{C2, as written, is internally impossible, and the corrected version
is sharper for it.} The post-ban basis cannot exist; the defensible objects are a
correctly-dated pre-ban basis trend and a price-discovery test on commodities that kept
trading, with the conditional-volatility question handled by a model that lets the data choose
the break date rather than imposing the ban date on it. Carried out, that ICSS test finds no
break at the suspension, and the original ``elevated volatility, significant at 10\%'' does not
reproduce --- on rebuilt data the cleanest commodity's GARCH is roughly flat, so the
conditional-volatility strand is left inconclusive and the volatility case is carried by C1.

\smallskip\noindent\textbf{C3 stays open by design} --- a team-led strand whose job is to turn
any market-level finding into a statement about consumers and farmers.

\smallskip\noindent The through-line for an interview is short and strong: we inherited three
results and, rather than dress them up, we verified them --- confirming none as stated,
overturning one, reframing another, and leaving the third honestly open. A project willing to
let its own rules kill its headline number is doing exactly what empirical research is meant
to do.

% =====================================================================
\part*{Part II --- Methodology Theory}
\addcontentsline{toc}{part}{Part II --- Methodology Theory}
\noindent\emph{Each method below is given as: the general theory, then why it fits this
problem, then how we implemented it. Scope is limited to the methods named in C1--C3.}

\section{The outcome variable: realized spot-price volatility}
\thy ``Volatility'' is the dispersion of price changes. Working from a daily price series
$P_t$, the daily log return is
\[
  r_t = \ln\!\left(\frac{P_t}{P_{t-1}}\right),
\]
and \emph{realized volatility} over a rolling window of $w$ trading days is the (sample)
standard deviation of those returns, annualised by $\sqrt{252}$:
\[
  \sigma_t \;=\; \sqrt{\frac{1}{w-1}\sum_{i=t-w+1}^{t}\big(r_i-\bar r\big)^2}\;\times\;\sqrt{252}.
\]
Log returns are used because they are (approximately) additive over time and symmetric in
percentage terms; annualisation puts the number on the familiar ``\% per year'' scale.

\why The suspension's \emph{stated} objective was to reduce price volatility, so volatility
is the natural dependent variable --- the policy must be judged on its own goal. A
rolling-window realized measure is model-free, transparent, and easy to defend: it makes no
distributional assumption and is directly interpretable, which matters when the headline
number (``+8--10\%'') is itself the claim.

\how We compute 30-day rolling realized volatility of daily log returns on the modal spot
("mandi") price, per geographic unit and commodity, annualised. A 60-day window is the
robustness twin. This $\sigma_t$ becomes the left-hand-side variable in the DiD/event-study
panel (C1) and motivates the conditional-variance model (C2).

\section{Difference-in-Differences (C1)}
\thy Difference-in-Differences (DiD) estimates a causal effect by comparing the
\emph{change} in an outcome for a treated group to the \emph{change} for an untreated
(control) group over the same period. The simplest $2\times2$ logic: if treated and control
would have moved in parallel absent treatment, then the treated group's extra movement after
treatment is the causal effect. In panel form with unit and time fixed effects,
\[
  y_{it} \;=\; \alpha_i \;+\; \lambda_t \;+\; \beta\,(\mathrm{Treat}_i \times \mathrm{Post}_t)
  \;+\; X_{it}'\gamma \;+\; \varepsilon_{it},
\]
where $\alpha_i$ are unit fixed effects (absorbing time-invariant differences across
commodities/districts), $\lambda_t$ are time fixed effects (absorbing common shocks in each
period), $\mathrm{Treat}_i$ marks banned commodities, $\mathrm{Post}_t$ marks dates after
20~Dec~2021, and $\beta$ --- the coefficient on their interaction --- is the DiD estimate.

\smallskip\noindent\textbf{Key assumptions.} (1) \emph{Parallel trends}: treated and control
would have followed the same trajectory absent the ban (the central, testable assumption).
(2) \emph{No anticipation}: behaviour does not change before the announcement. (3) \emph{No
spillovers / common-shock contamination}: shocks affecting both groups are absorbed by
$\lambda_t$; shocks hitting only one group threaten identification. Inference must account
for serial correlation within a commodity --- standard errors are clustered by commodity ---
and with few clusters a wild-cluster bootstrap is the honest correction.

\why The suspension is a textbook natural experiment: a sharp, dated policy that applied to
some commodities and not others. DiD is the canonical tool for exactly this structure. Its
great virtue here is that 2022 was full of \emph{common} shocks --- the Ukraine war, energy
prices, global food inflation --- and the control group differences these out: both treated
and control commodities faced the same macro environment, so a divergence between them is
attributable to the ban rather than the shock. This is the single strongest defence of the
design.

\how The panel is commodity$\times$geography$\times$month realized volatility. Treatment =
the banned set; controls = guar, castor, jeera, turmeric, cotton; the break is
20~Dec~2021. We estimate $\beta$ with unit and time fixed effects and commodity-clustered
standard errors. C1 is judged confirmed only if the sign matches, the magnitude lies in a
pre-registered $[4\%,16\%]$ band around the claimed 8--10\%, and it is significant at
conventional levels --- and only if the parallel-trends (pre-trend) test passes; if it
fails, the DiD is set aside in favour of a synthetic-control design.

\section{Event study (C1)}
\thy An event study is DiD with time resolution: instead of a single $\mathrm{Post}$ dummy,
the treatment interaction is split into a sequence of leads and lags around the event date,
\[
  y_{it} \;=\; \alpha_i + \lambda_t + \sum_{k\neq -1}\beta_k\,\big(\mathrm{Treat}_i \times \mathbf{1}[t-t^\ast=k]\big) + \varepsilon_{it},
\]
where $t^\ast$ is the event period and $k$ indexes periods before/after it (the period just
before, $k=-1$, is the omitted baseline). The pre-event coefficients ($k<0$) should be
statistically zero --- this is the formal \emph{test} of parallel trends --- while the
post-event coefficients ($k\ge 0$) trace the effect's \emph{dynamics}: immediate vs.\
delayed, transitory vs.\ persistent.

\why It does double duty: it validates the DiD's core assumption (flat pre-trends) and it
tells a richer story than a single average --- did volatility jump immediately at the
suspension and persist, or drift? For a policy debate that is exactly the question.

\how Estimated around 20~Dec~2021 on the same panel as the DiD, reporting the path of
$\beta_k$ with confidence bands; flat pre-period coefficients support C1, and the post-period
path characterises how the ban's volatility effect evolved.

\section{GARCH(1,1) (C2)}
\thy Financial and commodity returns exhibit \emph{volatility clustering} --- calm and
turbulent periods bunch together --- which a constant-variance model misses. GARCH
(Generalised Autoregressive Conditional Heteroskedasticity) lets today's variance depend on
recent shocks and recent variance. The GARCH(1,1) specification is
\[
  r_t = \mu + \varepsilon_t,\qquad \varepsilon_t = \sigma_t z_t,\quad z_t \sim \text{iid}(0,1),
\]
\[
  \sigma_t^2 \;=\; \omega \;+\; \alpha\,\varepsilon_{t-1}^2 \;+\; \beta\,\sigma_{t-1}^2 .
\]
Here $\alpha$ captures the reaction of variance to last period's shock, $\beta$ its
persistence; $\alpha+\beta$ (close to 1 in practice) measures how long volatility shocks
last, and stationarity requires $\alpha+\beta<1$. To test a policy effect, an exogenous
\emph{step dummy} for the post-ban regime is added to the variance equation,
\[
  \sigma_t^2 \;=\; \omega + \alpha\,\varepsilon_{t-1}^2 + \beta\,\sigma_{t-1}^2 + \gamma\,\mathrm{Post}_t,
\]
so $\gamma>0$ (and significant) means the suspension raised the \emph{conditional} variance
of the series, over and above the usual clustering dynamics.

\why GARCH complements DiD. DiD asks ``did average volatility shift across treated vs.\
control units?''; GARCH asks ``within a single commodity's price process, did the
time-series volatility regime change at the ban?''. The CV claim is specifically about a
GARCH(1,1) coefficient significant at 10\%, so reproducing that model is required to verify
or revise it. It is also the standard volatility-modelling tool, so an interviewer will
expect fluency in it.

\how A GARCH(1,1) is fit per commodity on the daily spot return series, with a post-ban step
dummy in the variance equation; the sign, size, and significance of $\gamma$ are the object
of interest (the ``significant at 10\%'' claim). EGARCH is the natural robustness extension
because it allows asymmetric (good-news/bad-news) volatility responses. \emph{Current verdict:}
the inherited ``elevated volatility, significant at 10\%'' does \emph{not} reproduce --- a
data-driven ICSS break-date test finds \emph{no} variance break near the suspension, and the
short-lived ``chana calmer'' reading proved to be a calendar-grid artifact (on clean data the
persistence is roughly flat). C2-volatility is therefore \emph{inconclusive}, and the
volatility verdict is carried by C1, not C2.

\section{The basis, and the C2 internal-inconsistency (C2)}
\thy The \emph{basis} is the difference between the futures price and the spot price of the
same commodity, $\text{basis}_t = F_t - S_t$. It reflects the cost of carry (storage,
financing) and convergence: as a contract nears expiry, futures and spot converge and the
basis shrinks toward zero. A ``widening'' basis means futures and spot are drifting apart ---
often read as a sign of stress or impaired arbitrage between the two markets.

\smallskip\noindent\textbf{The inconsistency.} C2 as originally stated claims a
\emph{``persistent basis widening post-ban''} for the banned commodities. But the basis
requires \emph{both} a futures and a spot price, and the suspension \emph{eliminated the
futures price} for exactly those commodities after 20~Dec~2021. A post-ban basis for a
banned commodity therefore cannot exist --- the claim is internally impossible as written.

\why Surfacing this is the point. It is a research-integrity check: a result that cannot
exist must be diagnosed before anything is rebuilt on top of it. The most likely
explanations are that the earlier work actually measured (a) a \emph{pre-ban} basis trend
and mislabelled its timing, (b) a \emph{control}-commodity basis (those still have futures),
or (c) a domestic--international price spread relabelled as ``basis''. Our reading of the
stale post-ban settlement quotes found in the raw MCX data (months of zero-volume marks)
suggests the mechanical trap that could have generated a spurious post-ban ``futures'' price.

\how We compute the basis only where it is defined: $F_t - S_t$ up to December 2021 for the
banned commodities (a pre-ban trend, correctly labelled), and across the full window for
control commodities (which retain futures). Any volatility claim that travels under the
``basis'' label is re-expressed in terms that are actually measurable, and the discrepancy
with the original statement is documented rather than papered over.

\section{C3 --- inflation transmission (overview)}
\thy\emph{(Overview only, by design.)} The transmission question asks whether suspending
futures fed through to the prices people actually pay and to farmers' decisions. Three
linked pieces: (i) \emph{price transmission} --- does the suspension move CPI/WPI food
sub-indices, tested with transmission regressions or vector-autoregression (VAR) style
dynamics between the policy, spot prices, and the inflation indices; (ii) \emph{acreage
response} --- do farmers plant less of a crop once the forward price signal from futures
disappears, in the spirit of Nerlove-type acreage models that relate planted area to
expected prices; (iii) \emph{district production} --- whether output responds, as a check on
the acreage channel.

\why It closes the loop from market microstructure to welfare: even if volatility changed,
the policy-relevant question is whether consumers and farmers were better or worse off. It is
the ``so what'' of the project.

\how\emph{(Conducted primarily by the wider team.)} The empirical objects are the CPI/WPI
food sub-indices (monthly), state/district acreage, and district production, related to the
suspension timing.

\paragraph{What the results would mean --- for each stakeholder and each case.}
\begin{center}
\small
\begin{tabular}{@{}p{2.4cm}p{3.5cm}p{3.5cm}p{3.5cm}@{}}
\toprule
 & \textbf{Ban ``worked''} (volatility/ inflation fell) & \textbf{Ban backfired} (volatility/ inflation rose) & \textbf{No effect} \\
\midrule
\textbf{Consumers} & Lower, steadier food prices --- the policy's intended win. & Paid more / more volatile prices --- policy counterproductive. & Prices driven by supply/demand and global factors, not futures. \\
\textbf{Farmers} & Mixed: steadier prices but lost a forward price signal and hedging tool. & Worse: lost hedging \emph{and} faced higher volatility. & Futures were not a binding channel for them. \\
\textbf{Government / SEBI} & Validates the intervention. & Argues for restoring futures; the scapegoating was misplaced. & Suspension was costly theatre --- no benefit, some cost. \\
\textbf{Exchanges / industry} & Hard to justify reopening. & Strong case to reopen; impaired price discovery hurt the ecosystem. & Neutral on prices, but lost business and liquidity. \\
\bottomrule
\end{tabular}
\end{center}
\noindent The honest framing in interview: C3 is an active, team-led strand; the value of
the evaluation is that it is two-armed --- it weighs any stability benefit against the cost
to price discovery and hedging, rather than advocating a position.

\end{document}
